The migration and modernization of legacy robotic systems involve adapting them to the mechanisms of new frameworks and refactoring their implementation while seeking to preserve the essential behaviors of the original system. In this context, this work proposes a systematic migration approach guided by the establishment of a behavioral \textit{ground truth} in the legacy system. The approach integrates behavioral specification and verification, the investigation of how these behaviors are implemented in the legacy system, and the analysis of transformations carried out during migration. The approach was examined through a case study of the migration of the \textit{Body Sensor Network} (BSN) from ROS 1 to ROS 2, combining the retrospective reconstruction of activities with complementary analyses of the implementations. The results indicate that the BSN migration process can be retrospectively systematized by establishing relationships among behavioral contracts, implementation responsibilities, and transformations identified between ROS 1 and ROS 2. The analysis made it possible to characterize changes in these responsibilities and identify evidence of behavioral preservation by comparing BDD scenarios, their test oracles, and execution results across both versions. The study also identified limitations in reconstructing the original migration process and in the scope of the available evidence of behavioral preservation.